% !TEX program = lualatex
\documentclass[aspectratio=43,professionalfonts,handout]{beamer}
\usepackage{beamer_template}
\usepackage{enumerate}

\newcommand{\LectureNum}{11}
\newcommand{\LectureTitle}{第3章 フーリエ解析　2節 フーリエ変換と積分定理}
\newcommand{\TextPages}{p.\,91-104}
\newcommand{\CourseName}{応用数学}
\renewcommand{\TermName}{前期}

\begin{document}

\begin{frame}{第3章 フーリエ解析　2節 フーリエ変換と積分定理}
  \begin{block}{本節の内容}
  \begin{itemize}
    \item フーリエ積分定理とフーリエ変換　（p.\,91--95）
    \item フーリエ変換の性質と公式　（p.\,96--99）
    \item スペクトル　（p.\,100--103）
  \end{itemize}
  \end{block}
  \vfill\hfill{\scriptsize \textcolor{gray}{教科書 p.\,91--104}}
\end{frame}

\begin{frame}{フーリエ積分定理とフーリエ変換}
  \begin{block}{}
  {\large フーリエ積分定理とフーリエ変換}
  \end{block}
  \vfill\hfill{\scriptsize \textcolor{gray}{p.\,91--95}}
\end{frame}

\begin{frame}{定義：フーリエ変換}
  \begin{block}{}
  \end{block}
  \vfill\hfill{\scriptsize \textcolor{gray}{p.\,91--95}}
\end{frame}

\begin{frame}{定理：フーリエ積分定理}
  \begin{block}{}
  \textbf{条件}：
  \begin{itemize}
    \item $f(t)$ は任意の有限区間で区分的に滑らか
    \item $\int_{-\infty}^{\infty} |f(t)| dt < \infty$（絶対可積分）
  \end{itemize}
  \[
    f(t) = (1/\pi) \int_{0}^\infty [\int_{-\infty}^{\infty} f(\tau)\cos(\omega(t-\tau)) d\tau] d\omega

  \]
  \end{block}
  \vfill\hfill{\scriptsize \textcolor{gray}{p.\,91--95}}
\end{frame}

\begin{frame}{フーリエ余弦変換}
  \begin{block}{}
  \[
    Fc(\omega) = \int_{0}^\infty f(t)\cos(\omega t) dt
  \]
  \[
    f(t) = (2/\pi) \int_{0}^\infty Fc(\omega)\cos(\omega t) d\omega
  \]
  {\footnotesize $f(t)$ が偶関数の場合に対応}
  \end{block}
\end{frame}

\begin{frame}{フーリエサイン変換}
  \begin{block}{}
  \[
    Fs(\omega) = \int_{0}^\infty f(t)\sin(\omega t) dt
  \]
  \[
    f(t) = (2/\pi) \int_{0}^\infty Fs(\omega)\sin(\omega t) d\omega
  \]
  {\footnotesize $f(t)$ が奇関数の場合に対応}
  \end{block}
\end{frame}

\begin{frame}{例題1　（p.\,91）}
  \begin{exampleblock}{問題}
    $f(x) = e^{-|x|}$ のフーリエ変換を求めよ\\[2pt]
  \end{exampleblock}
\end{frame}

\begin{frame}{例題1　解答}
  \begin{exampleblock}{解答}
  \begin{enumerate}[1.]
    \item[] $F(u) = \int_{-\infty}^{\infty} e^{-|x|} e^{-iux} dx$
    \item[] $= \int_{-\infty}^{0} e^{x} e^{-iux} dx + \int_{0}^{\infty} e^{-x} e^{-iux} dx$
    \item[] $= 1/(1-iu) + 1/(1+iu)$
  \end{enumerate}
  \end{exampleblock}
  \begin{alertblock}{答}
    $F(u) = 2/(1+u^{2})$
  \end{alertblock}
  {\footnotesize \textcolor{gray}{\textit{$lim_{x \to ±\infty} e^{(1-iu)x} = 0$ を用いる}}}
\end{frame}

\begin{frame}{例題2　（p.\,92）}
  \begin{exampleblock}{問題}
    $f(x) = \begin{cases} 1 & (|x| \leq 1) \\ 0 & (|x| > 1) \end{cases}$ のフーリエ変換を求めよ\\[2pt]
  \end{exampleblock}
\end{frame}

\begin{frame}{例題2　解答}
  \begin{exampleblock}{解答}
  \begin{enumerate}[1.]
    \item[] $F(u) = \int_{-1}^{1} e^{-iux} dx = [-e^{-iux}/(iu)]_{-1}^{1}$
    \item[] $= (e^{iu} - e^{-iu})/(iu) = 2sin u / u$
  \end{enumerate}
  \end{exampleblock}
  \begin{alertblock}{答}
    $F(u) = 2sin u / u$
  \end{alertblock}
\end{frame}

\begin{frame}{例題3　（p.\,95）}
  \begin{exampleblock}{問題}
    $f(x) = \begin{cases} 1-|x| & (|x| \leq 1) \\ 0 & (|x| > 1) \end{cases}$ のフーリエ余弦変換を求めよ\\[2pt]
  \end{exampleblock}
\end{frame}

\begin{frame}{例題3　解答}
  \begin{exampleblock}{解答}
  \begin{enumerate}[1.]
    \item[] $f(x)$ は偶関数なのでフーリエ余弦変換を用いる
    \item[] $F(u) = 2\int_{0}^{1} (1-x) \cos(ux) dx$
    \item[] $= 2\begin{cases} [(1-x)\sin(ux)/u]_{0}^1 + (1/u)\int_{0}^1 \sin(ux) dx \end{cases}$
    \item[] $= 2\begin{cases} 0 + (1/u)[-\cos(ux)/u]_{0}^1 \end{cases} = 2(1-\cos u)/u^{2}$
  \end{enumerate}
  \end{exampleblock}
  \begin{alertblock}{答}
    $F(u) = 2(1-\cos u)/u^{2}$
  \end{alertblock}
\end{frame}

\begin{frame}{演習問題（p.\,91--95）}
  \begin{exampleblock}{自分で解いてみよう}
  \begin{description}
    \item[問1] 次の関数のフーリエ変換を求めよ
    $(1) f(x) = { 0  (x > 0),  e^x  (x \leq 0) }$\\
    $(2) g(x) = { xe^{-x}  (x > 0),  0  (x \leq 0) }$\\
    \item[問2] 次の関数のフーリエ変換を求めよ
    $(1) f(x) = { 1  (0 \leq x \leq 1),  0  (x < 0, x > 1) }$\\
    $(2) g(x) = { x  (|x| \leq 1),  0  (|x| > 1) }$\\
    \item[問3] 問2(1)の関数に積分定理を適用して、次の等式を証明せよ
    $(1) (1/2\pi)\int_{-\infty}^{\infty} (1-e^{-iu})/iu \cdot e^{iux} du = { 1  (0<x<1),  1/2  (x=0,1),  0  (x<0, x>1) }$\\
    $(2) \int_{-\infty}^{\infty} \sin u / u  du = \pi$\\
    \item[問4] 次の関数のフーリエサイン変換を求めよ
    $f(x) = { -1  (-1 \leq x < 0),  1  (0 \leq x \leq 1),  0  (|x| > 1) }$\\
  \end{description}
  \end{exampleblock}
\end{frame}

\begin{frame}{フーリエ変換の性質と公式}
  \begin{block}{}
  {\large フーリエ変換の性質と公式}
  \end{block}
  \vfill\hfill{\scriptsize \textcolor{gray}{p.\,96--99}}
\end{frame}

\begin{frame}{フーリエ変換の性質}
  \begin{block}{}
  \begin{description}
    \item[I. 線形性]\mbox{}\\$F[c_{1}f_{1}(x) + c_{2}f_{2}(x)] = c_{1}F_{1}(u) + c_{2}F_{2}(u)  (c_{1}, c_{2} は定数)$
    \item[II. 移動則（時間シフト）]\mbox{}\\$F[f(x-a)] = e^{-iau} F(u)  (a は実数)$
    \item[III. 移動則（周波数シフト）]\mbox{}\\$F[e^{iax}f(x)] = F(u-a)  (a は実数)$
    \item[IV. 相似性（スケーリング）]\mbox{}\\$F[f(ax)] = (1/|a|) F(u/a)  (a は 0 でない実数)$
    \item[V. 微分則]\mbox{}\\$F[f^{(n)}(x)] = (iu)^n F(u)  (n は自然数)$
    \item[VI. 双対微分則]\mbox{}\\$F[(-ix)^n f(x)] = F^{(n)}(u)  (n は自然数)$
  \end{description}
  \end{block}
  \vfill\hfill{\scriptsize \textcolor{gray}{p.\,96--99}}
\end{frame}

\begin{frame}{パーセバルの等式}
  \begin{block}{}
  \[
    \int_{-\infty}^\infty |f(t)|^{2} dt = (1/2\pi) \int_{-\infty}^\infty |F(\omega)|^{2} d\omega
  \]
  \end{block}
\end{frame}

\begin{frame}{例題4　（p.\,99）}
  \begin{exampleblock}{問題}
    次の等式を証明せよ\\[2pt]
    $(1) F[e^{-a^{2}x^{2}} * e^{-b^{2}x^{2}}] = (\pi/(ab)) e^{-u^{2}(a^{2}+b^{2})/(4a^{2}b^{2})}$ 等（たたみこみのフーリエ変換と 98 ページの公式 (1) を利用）\\[2pt]
    (2) 上記の結果と比較して、畳み込み関数の閉じた形を求める\\[2pt]
  \end{exampleblock}
\end{frame}

\begin{frame}{例題4　解答}
  \begin{exampleblock}{解答}
  \end{exampleblock}
  {\footnotesize \textcolor{gray}{\textit{たたみこみ定理 $: F[f*g] = F(u)G(u)$ を適用}}}
\end{frame}

\begin{frame}{演習問題（p.\,96--99）}
  \begin{exampleblock}{自分で解いてみよう}
  \begin{description}
    \item[問5] 性質 III および VI を証明せよ
    \item[問6] $F[f(x)] = F(u), F[g(x)] = G(u)$ のとき、次の等式を証明せよ
    $F[f(x)g(x)] = (1/(2\pi)) F(u) * G(u)$\\
    \item[問7] 96ページの性質を用いて、次の関数のフーリエ変換を求めよ
    $(1) xe^{-x^{2}/2}$\\
    $(2) x^{2}e^{-x^{2}/2}$\\
    \item[問8] 次の等式を証明せよ
    $F^{-}^{1}[e^{-bu^{2}}] = 1/(2\sqrt\pi b) \cdot e^{-x^{2}/(4b)}$（b は正の定数）\\
    \item[問9] 次の等式を証明せよ
    $(1) F[e^{-x^{2}/2} * xe^{-x^{2}/2}] = -2\pi iu\cdot e^{-u^{2}}$\\
    $(2) e^{-x^{2}/2} * xe^{-x^{2}/2} = (\sqrt\pi/2)\cdot xe^{-x^{2}/4}$\\
  \end{description}
  \end{exampleblock}
\end{frame}

\begin{frame}{スペクトル}
  \begin{block}{}
  {\large スペクトル}
  \end{block}
  \vfill\hfill{\scriptsize \textcolor{gray}{p.\,100--103}}
\end{frame}

\begin{frame}{例2　（p.\,100）}
  \begin{exampleblock}{問題}
    周期4の矩形波のスペクトルを求めよ\\[2pt]
    $f(x) = \begin{cases} 1 & (-1 < x < 1) \\ 1/2 & (x = 1, 3) \\ 0 & (1 < x < 3) \end{cases},  f(x+4) = f(x)$\\[2pt]
  \end{exampleblock}
\end{frame}

\begin{frame}{例2　解答}
  \begin{exampleblock}{解答}
  \end{exampleblock}
  {\footnotesize \textcolor{gray}{\textit{線スペクトル（ $\omega$ の値がとびとびに現れる）}}}
\end{frame}

\begin{frame}{例3　（p.\,101）}
  \begin{exampleblock}{問題}
    次の関数のスペクトルを求めよ\\[2pt]
    $f(x) = \begin{cases} 1 & (|x| < 1) \\ 1/2 & (|x| = 1) \\ 0 & (|x| > 1) \end{cases}$\\[2pt]
  \end{exampleblock}
\end{frame}

\begin{frame}{例3　解答}
  \begin{exampleblock}{解答}
  \end{exampleblock}
  {\footnotesize \textcolor{gray}{\textit{連続スペクトル（ $S_i(\omega)$ の値が連続的に変化する）}}}
\end{frame}

\begin{frame}{演習問題（p.\,100--103）}
  \begin{exampleblock}{自分で解いてみよう}
  \begin{description}
    \item[問10] 次の関数のスペクトルを求めよ
    $f(x) = 1 - |x|  (|x| \leq 1),  f(x+2) = f(x)$\\
    \item[問11] 次の関数のスペクトルを求めよ
    $f(x) = { 1-|x|  (|x| \leq 1),  0  (|x| > 1) }$\\
  \end{description}
  \end{exampleblock}
\end{frame}

\begin{frame}{練習問題2　（p.\,104）}
  \begin{block}{}
  \begin{description}
    \item[問1] 次の関数のフーリエ変換を求めよ
    $(1) f(x) = { 2  (0 \leq x < 2),  1  (2 \leq x < 3),  0  (x < 0, x \geq 3) }$\\
    $(2) g(x) = { 1 - x/2  (0 \leq x < 2),  0  (x < 0, x \geq 2) }$\\
    \item[問2] 次の関数について、(1) はフーリエ正弦変換、(2) はフーリエ余弦変換をそれぞれ求めよ
    $(1) f(x) = { \sin x  (|x| \leq \pi),  0  (|x| > \pi) }$\\
    $(2) g(x) = { \cos x  (|x| \leq \pi/2),  0  (|x| > \pi/2) }$\\
  \end{description}
  \end{block}
\end{frame}

\end{document}
